\section{Dimension Theory}

\begin{exercise}\label{ex11.1}
	Move $P$ to the origin to remove the constant term of $f$. $\mathfrak{m}=\left(x_1,\cdots,x_n\right)/(f)$, so $\mathfrak{m}/\mathfrak{m}^2=\left.\left(x_1,\cdots,x_n\right)\middle/\left(\left(x_1,\cdots,x_n\right)^2+(f)\right)\right.$. $\left\{x_i\right\}_{i=1}^n$ is a set of generators $\mathfrak{m}/\mathfrak{m}^2$. If $f\in\left(x_1,\cdots,x_n\right)^2$, $\left\{x_i\right\}_{i=1}^n$ is a basis. If $f\notin\left(x_1,\cdots,x_n\right)^2$, the existence of linear term of $f$ means $\left\{x_i\right\}_{i=1}^n$ is linearly dependent.
	
	Use \Cref{11.18} on $A_\mathfrak{m}=k\left[x_1,\cdots,x_n\right]_\mathfrak{m}/(f)_\mathfrak{m}$, then \Cref{11.22}.
\end{exercise}

\begin{exercise}\label{ex11.2}
	
\end{exercise}

\begin{exercise}\label{ex11.3}

\end{exercise}

\begin{exercise}\label{ex11.4}

\end{exercise}

\begin{exercise}\label{ex11.5}

\end{exercise}

\begin{exercise}\label{ex11.6}
	From \hyperref[ex4.7]{\Cref*{ex4.7}(ii)}, if $\mathfrak{p}_0\subsetneq\cdots\subsetneq\mathfrak{p}_n$ is a chain in $A$, then $\mathfrak{p}_0[x]\subsetneq\cdots\subsetneq\mathfrak{p}_n[x]\subsetneq\mathfrak{p}_n[x]+(x)$ is a chain in $A[x]$.
	
	Let $f:A\rightarrow A[x]$ be the embedding. The fiber of $f^*$ at $\mathfrak{p}$ is $\operatorname{Spec}\left(k(\mathfrak{p})\otimes_AA[x]\right)=\operatorname{Spec}(k(\mathfrak{p})[x])$. Since $\dim\left(k(\mathfrak{p})[x]\right)=1$, a chain in $A[x]$ whose contraction is $\mathfrak{p}$ contains no more than 2 primes, so a chain in $A[x]$ contains at most $2+2\dim A$ prime ideals.
\end{exercise}

\begin{exercise}\label{ex11.7}
	Localize at $\mathfrak{p}$, we denote $m=\operatorname{height}\mathfrak{p}=\dim A_\mathfrak{p}$. By \Cref{11.14}, there is a $\mathfrak{p}$-primary ideal $\mathfrak{q}$ generated by $m$ elements, so $\mathfrak{q}[x]$ also generated by $m$ elements. From \hyperref[ex4.7]{\Cref*{ex4.7}(iii)}, $\mathfrak{q}[x]$ is $\mathfrak{p}[x]$-primary, then use \Cref{11.14} again we have $\operatorname{height}\mathfrak{p}[x]\leqslant m$. Then use \Cref{ex11.6}.
\end{exercise}